% SPDX-License-Identifier: Apache-2.0
\section{Signals, Ports and Interfaces}
\label{sec:e-wiring}

These three were the muddiest part of the merged specification, and LHDL's
version of them is where its grammar and its examples disagreed most. The
embedding settles them, and the settlement is one idea.

\begin{rulebox}[Direction is not a property of a wire]
The same wire is driven at one end and read at the other. Direction
belongs to the end, not to the wire, so a unit never holds a wire. It
holds an end.
\end{rulebox}

\subsection{A signal has two ends}

Creating a signal yields its ends separately, the way
\code{std::sync::mpsc} yields a sender and a receiver.

\begin{lstlisting}[caption={A signal is created as a pair. \code{T} is a struct, so a wire may be a bundle without a second concept for the bundled case.},label={lst:split}]
pub fn signal<T: Copy + Default>() -> (Out<T>, In<T>);

let (tx, rx) = signal::<Beat>();
\end{lstlisting}

\code{Out<T>} has \code{set} and no \code{get}. \code{In<T>} has \code{get}
and no \code{set}. A unit given the reading end cannot drive, and the
compiler says so in the plainest possible terms:
\code{error[E0599]: no method named `set` found for struct `In`}.

\subsection{Cardinality is the opposite of mpsc, and the types say so}

A channel in software has many senders and one receiver. A wire has one
driver and many readers. The library inverts the derive accordingly.

\begin{table}[H]
\caption{Why \code{In} is \code{Clone} and \code{Out} is not.}
\label{tab:cardinality}
\centering
\footnotesize
\begin{tabular}{@{}lll@{}}
\toprule
 & Sender or driver & Receiver or reader \\
\midrule
\code{mpsc} & \code{Clone}     & unique \\
a wire      & unique           & \code{Clone} \\
\bottomrule
\end{tabular}
\end{table}

That one line of derive buys the oldest error message in hardware design.
A second driver on a net would be a second \code{Out}, \code{signal}
returns exactly one, and \code{Out} is not \code{Clone}, so there is no
way to make another.

\begin{lstlisting}[caption={Two drivers on one wire. The classic multiple-driver error is a move error.},label={lst:twodrivers}]
let (tx, _rx) = signal::<u32>();
let a = Driver { out: tx };
let b = Driver { out: tx };
// error[E0382]: use of moved value: `tx`
\end{lstlisting}

Nothing in the language had to define that rule. Fanout, meanwhile, is
\code{rx.clone()} and costs nothing, because reading a wire is free.

\subsection{An interface is a struct of ends}

An interface is a named group. Its members are signals and channels, and
it holds no direction, because its members do not.

A \term{role} is the pairing that says which end of each member this side
takes. A \term{port} is what a unit is handed: an interface's ends,
selected by a role.

\begin{lstlisting}[caption={An interface and the two roles over it. Each role names the end it takes for every member.},label={lst:iface}]
pub struct Wishbone {
    pub adr: Signal<U<32>>,
    pub ack: Signal<Bit>,
    pub dat: Chan<Beat>,
}

pub struct InitiatorPort {
    pub adr: Out<U<32>>,
    pub ack: In<Bit>,
    pub dat: Tx<Beat>,
}

pub struct TargetPort {
    pub adr: In<U<32>>,
    pub ack: Out<Bit>,
    pub dat: Rx<Beat>,
}
\end{lstlisting}

The two role structs are mechanical: each field is the same member with
the other end. A derive writes them, which is what the \code{modport}
keyword used to do, and it cannot get them out of step with the interface
because it reads the interface to make them.

\subsection{Where the ends go}

A unit does not store a reference to the interface. It is handed its port
when \code{run} is called, which is what the inputs and outputs of
\code{Module} are for, and what makes wiring a call rather than a graph of
back-references.

\begin{lstlisting}[caption={The parent creates the interface and hands each child the end it takes.},label={lst:wire-up}]
impl Module<(), ()> for Top {
    async fn run(&mut self, _i: (), _o: ()) {
        join2(
            self.cpu.run(self.bus.initiator(), ()),
            self.ram.run(self.bus.target(), ()),
        ).await;
    }
}
\end{lstlisting}

\begin{table*}[t]
\caption{The four concepts, and what each one is responsible for. Reading
down the table is the life of a wire.}
\label{tab:wiring}
\centering
\footnotesize
\begin{tabular}{@{}llll@{}}
\toprule
Concept & Is & Knows about direction & Held by \\
\midrule
\code{Signal<T>}, \code{Chan<T>} & one wire, or one handshaked channel & no & nobody \\
\code{Out<T>} / \code{In<T>}, \code{Tx<T>} / \code{Rx<T>} & one end of one member & yes, by which methods exist & a unit \\
interface & a named group of members & no & the unit containing both ends \\
role and port & a group of ends, one per member & yes, by construction & passed to \code{run} \\
\bottomrule
\end{tabular}
\end{table*}
